
\section*{3. Coalgebras and Cartier Duality}
\phantomsection
\label{sga3:VIIA:3}
\addcontentsline{toc}{section}{3. Coalgebras and Cartier Duality}

\medskip
\noindent\textbf{3.1.}
\phantomsection
\label{sga3:VIIA:3.1}
Let \(S\) be a scheme (or, more generally, a ringed space). An
\(\mathcal O_S\)-\emph{coalgebra}\footnote{The term \emph{cogebra} is
also used; see Bourbaki, \emph{Algebra}, III, \S 11.1. Notice also
that in this expos\'e, as well as in Expos\'e VII\(_B\), we work in
the category of cocommutative coalgebras (that is, coalgebras
satisfying condition~\SGARefLink{sga3:vii:paragraph:5-1:equation:i}{(i)}). This is crucial for defining the product
and the notion of a coalgebra in groups; see \SGARefLink{sga3:VIIA:3.1.0}{3.1.0} and \SGARefLink{sga3:VIIA:3.2}{3.2}.}
is a pair \((\mathcal U,\Delta_{\mathcal U})\) consisting of an
\(\mathcal O_S\)-module \(\mathcal U\) and a morphism of
\(\mathcal O_S\)-modules
\[
  \Delta_{\mathcal U}:\mathcal U\longrightarrow
  \mathcal U\otimes_{\mathcal O_S}\mathcal U,
\]
called the \emph{diagonal morphism} or \emph{comultiplication}, such
that:

\smallskip
\noindent
\SGARefLink{sga3:vii:paragraph:5-1:equation:i}{(i)} \(\sigma\circ\Delta_{\mathcal U}=\Delta_{\mathcal U}\), where
\(\sigma(a\otimes b)=b\otimes a\).

\smallskip
\noindent
\SGARefLink{sga3:vii:paragraph:5-1:equation:ii}{(ii)} The square
\begin{center}
  \phantomsection
  \label{sga3:VIIA:3.1-coassociativity-square}
  \begin{tikzcd}[column sep=huge,row sep=large]
\mathcal U
  \arrow[r,"\Delta_{\mathcal U}"]
  \arrow[d,"\Delta_{\mathcal U}"']
&
\mathcal U\otimes_{\mathcal O_S}\mathcal U
  \arrow[d,"\operatorname{id}_{\mathcal U}\otimes\Delta_{\mathcal U}"]
\\
\mathcal U\otimes_{\mathcal O_S}\mathcal U
  \arrow[r,"\Delta_{\mathcal U}\otimes\operatorname{id}_{\mathcal U}"]
&
\mathcal U\otimes_{\mathcal O_S}\mathcal U\otimes_{\mathcal O_S}\mathcal U
\end{tikzcd}

\end{center}
is commutative.

\smallskip
\noindent
(iii) There exists a morphism of \(\mathcal O_S\)-modules
\[
  \varepsilon_{\mathcal U}:\mathcal U\longrightarrow\mathcal O_S,
\]
called the \emph{augmentation}, such that the two composites
\begin{center}
  \phantomsection
  \label{sga3:VIIA:3.1-augmentation-sequences}
  \(
\begin{gathered}
\mathcal U
  \xrightarrow{\Delta_{\mathcal U}}
\mathcal U\otimes_{\mathcal O_S}\mathcal U
  \xrightarrow{\operatorname{id}_{\mathcal U}\otimes\varepsilon_{\mathcal U}}
\mathcal U\otimes_{\mathcal O_S}\mathcal O_S
\simeq\mathcal U
\\[0.6em]
\mathcal U
  \xrightarrow{\Delta_{\mathcal U}}
\mathcal U\otimes_{\mathcal O_S}\mathcal U
  \xrightarrow{\varepsilon_{\mathcal U}\otimes\operatorname{id}_{\mathcal U}}
\mathcal O_S\otimes_{\mathcal O_S}\mathcal U
\simeq\mathcal U
\end{gathered}
\)

\end{center}
are the identity morphism of \(\mathcal U\).

If \(\varepsilon_{\mathcal U}\) and
\(\varepsilon'_{\mathcal U}\) are two augmentations, then
\[
  \varepsilon_{\mathcal U}
  =(\varepsilon_{\mathcal U}\otimes\varepsilon'_{\mathcal U})
    \circ\Delta_{\mathcal U}
  =\varepsilon'_{\mathcal U};
\]
the augmentation is therefore uniquely determined by~(iii).

If \((\mathcal U,\Delta_{\mathcal U})\) and
\((\mathcal V,\Delta_{\mathcal V})\) are two
\(\mathcal O_S\)-coalgebras, a \emph{morphism} from the first to the
second is a morphism of \(\mathcal O_S\)-modules
\(f:\mathcal U\to\mathcal V\) for which the diagrams
\begin{center}
  \phantomsection
  \label{sga3:VIIA:3.1-coalgebra-morphism-diagrams}
  \begin{minipage}{0.47\linewidth}
\centering
\begin{tikzcd}[column sep=large,row sep=large]
\mathcal U
  \arrow[r,"f"]
  \arrow[d,"\Delta_{\mathcal U}"]
&
\mathcal V
  \arrow[d,"\Delta_{\mathcal V}"]
\\
\mathcal U\otimes\mathcal U
  \arrow[r,"f\otimes f"]
&
\mathcal V\otimes\mathcal V
\end{tikzcd}
\end{minipage}
\hfill
\text{and}
\hfill
\begin{minipage}{0.37\linewidth}
\centering
\begin{tikzcd}[column sep=large,row sep=large]
\mathcal U
  \arrow[rr,"f"]
  \arrow[dr,"\varepsilon_{\mathcal U}"']
&&
\mathcal V
  \arrow[dl,"\varepsilon_{\mathcal V}"]
\\
&
\mathcal O_S
\end{tikzcd}
\end{minipage}

\end{center}
are commutative. Morphisms of coalgebras compose like morphisms of
\(\mathcal O_S\)-modules, so we may speak of the
\emph{category of \(\mathcal O_S\)-coalgebras}.

\medskip
\noindent\textbf{3.1.0.}
\phantomsection
\label{sga3:VIIA:3.1.0}
\footnote{The editors added the numbering \SGARefLink{sga3:VIIA:3.1.0}{3.1.0} for later
references.}
This category has finite products: its final object is the
\(\mathcal O_S\)-module \(\mathcal O_S\), with the identity as
comultiplication. The product of two coalgebras
\((\mathcal U,\Delta_{\mathcal U})\) and
\((\mathcal V,\Delta_{\mathcal V})\) is the tensor product
\(\mathcal U\otimes_{\mathcal O_S}\mathcal V\), whose
comultiplication is the composite
\begin{center}
  \phantomsection
  \label{sga3:VIIA:3.1.0-product-comultiplication}
  \resizebox{0.98\linewidth}{!}{%
\(
\mathcal U\otimes\mathcal V
\xrightarrow{\Delta_{\mathcal U}\otimes\Delta_{\mathcal V}}
\mathcal U\otimes\mathcal U\otimes\mathcal V\otimes\mathcal V
\xrightarrow{\operatorname{id}_{\mathcal U}\otimes\sigma\otimes\operatorname{id}_{\mathcal V}}
\mathcal U\otimes\mathcal V\otimes\mathcal U\otimes\mathcal V
\)
}

\end{center}
where \(\sigma(a\otimes b)=b\otimes a\). The canonical projections
of \(\mathcal U\otimes\mathcal V\) onto the factors
\(\mathcal U\) and \(\mathcal V\) are
\(\operatorname{id}_{\mathcal U}\otimes\varepsilon_{\mathcal V}\)
and
\(\varepsilon_{\mathcal U}\otimes\operatorname{id}_{\mathcal V}\),
respectively.\footnote{The preceding assertion was added by the
editors. Recall also that, to show that
\(\mathcal U\otimes\mathcal V\) is indeed the product of
\(\mathcal U\) and \(\mathcal V\) in the category of cocommutative
\(\mathcal O_S\)-cogebras, one argues as follows. Given an arbitrary
\(\mathcal O_S\)-cogebra \(\mathcal E\) and cogebra morphisms
\(f:\mathcal E\to\mathcal U\) and
\(g:\mathcal E\to\mathcal V\), any cogebra morphism
\(\varphi:\mathcal E\to\mathcal U\otimes\mathcal V\) satisfying
\(\operatorname{pr}_{\mathcal U}\circ\varphi=f\) and
\(\operatorname{pr}_{\mathcal V}\circ\varphi=g\) must equal
\((f\otimes g)\circ\Delta_{\mathcal E}\); and this is a cogebra
morphism if and only if it also equals
\((g\otimes f)\circ\Delta_{\mathcal E}\).}
The ``diagonal morphism''
\(\mathcal U\to\mathcal U\otimes\mathcal U\), corresponding to the
pair
\((\operatorname{id}_{\mathcal U},
  \operatorname{id}_{\mathcal U})\),
is precisely the comultiplication \(\Delta_{\mathcal U}\).

\medskip
\noindent\textbf{3.1.1.}
\phantomsection
\label{sga3:VIIA:3.1.1}
Let \(\mathcal A\) be a commutative \(\mathcal O_S\)-algebra which,
as an \(\mathcal O_S\)-module, is locally free and of finite type.
Put
\[
  \mathcal A^{*}
  =\mathcal Hom_{\mathcal O_S\text{-Mod.}}
    (\mathcal A,\mathcal O_S).
\]
The canonical morphism
\[
  \varphi:\mathcal A^{*}\otimes_{\mathcal O_S}\mathcal A^{*}
  \longrightarrow
  (\mathcal A\otimes_{\mathcal O_S}\mathcal A)^{*}
\]
is invertible. If
\(m:\mathcal A\otimes\mathcal A\to\mathcal A\) is the morphism
defining the multiplication of \(\mathcal A\), composition gives a
diagonal morphism
\begin{center}
  \phantomsection
  \label{sga3:VIIA:3.1.1-dual-coalgebra-sequence}
  \(
\Delta_{\mathcal A^*}:\quad
\mathcal A^*
\xrightarrow{m^*}
(\mathcal A\otimes\mathcal A)^*
\xrightarrow{\varphi^{-1}}
\mathcal A^*\otimes\mathcal A^*\,.
\)

\end{center}
This diagonal morphism plainly makes \(\mathcal A^{*}\) into an
\(\mathcal O_S\)-coalgebra whose augmentation is the transpose of
the morphism \(\mathcal O_S\to\mathcal A\) defined by the unit
section of \(\mathcal A\). Moreover, it is clear that:

\begin{quote}
\itshape
The functor \(\mathcal A\mapsto\mathcal A^{*}\) is an
anti-equivalence from the category of \(\mathcal O_S\)-algebras
which are locally free and of finite type as
\(\mathcal O_S\)-modules onto the category of
\(\mathcal O_S\)-coalgebras which are locally free and of finite
type as \(\mathcal O_S\)-modules.
\end{quote}

\medskip
\noindent\textbf{3.1.2.}
\phantomsection
\label{sga3:VIIA:3.1.2}
To every \(\mathcal O_S\)-coalgebra \(\mathcal U\) there is
canonically associated an \(S\)-functor
\begin{center}
  \phantomsection
  \label{sga3:VIIA:3.1.2-spec-star-functor}
  \(
\operatorname{Spec}^{*}\mathcal U:
(\mathbf{Sch}/S)^{\circ}\longrightarrow(\mathbf{Ens})
\,.
\)

\end{center}
Indeed, for every \(S\)-scheme \(q:T\to S\),
\[
  q^{*}(\mathcal U\otimes_{\mathcal O_S}\mathcal U)
  \simeq
  q^{*}(\mathcal U)\otimes_{\mathcal O_T}q^{*}(\mathcal U),
\]
so \(q^{*}(\Delta_{\mathcal U})\) makes
\(\mathcal U_T=q^{*}(\mathcal U)\) into an
\(\mathcal O_T\)-coalgebra. Thus, by definition and with an obvious
abuse of notation,\footnote{For
\(x\otimes y\in
 \Gamma(T,\mathcal U_T)\otimes_{\mathcal O(T)}
 \Gamma(T,\mathcal U_T)\),
its image in
\(\Gamma(T,\mathcal U_T\otimes_{\mathcal O_T}\mathcal U_T)\)
is again denoted \(x\otimes y\).}
\[
  (\operatorname{Spec}^{*}\mathcal U)(T)
  =
  \left\{
    x\in\Gamma(T,\mathcal U_T)
    \ \middle|\
    \varepsilon_{\mathcal U_T}(x)=1,\quad
    \Delta_{\mathcal U_T}(x)=x\otimes x
  \right\}.
\]

Sections \(x\) of \(\mathcal U_T\) plainly correspond to morphisms
of \(\mathcal O_T\)-modules
\(\xi:\mathcal O_T\to\mathcal U_T\). The conditions
\(\varepsilon(x)=1\) and \(\Delta(x)=x\otimes x\) say exactly that
\(\xi\) is a morphism of coalgebras. Hence also
\[
  (\operatorname{Spec}^{*}\mathcal U)(T)
  =
  \operatorname{Hom}_{\mathcal O_T\text{-coalg.}}
    (\mathcal O_T,\mathcal U_T).
\]
In particular:

\medskip
\noindent\textbf{Proposition 3.1.2.1.}
\phantomsection
\label{sga3:VIIA:3.1.2.1}
\footnote{The editors added the numbering \SGARefLink{sga3:VIIA:3.1.2.1}{3.1.2.1} for later
references. Notice also that the \(S\)-functor
\(\operatorname{Spec}^{*}\mathcal U\) is a sheaf for the Zariski
topology; it is even an fpqc sheaf if \(\mathcal U\) is a
quasi-coherent \(\mathcal O_S\)-module.}
\textit{Let \(\mathcal A\) be a commutative
\(\mathcal O_S\)-algebra which is locally free of finite type as an
\(\mathcal O_S\)-module. Then the \(S\)-functor
\(\operatorname{Spec}^{*}\mathcal A^{*}\) is represented by
\(\operatorname{Spec}\mathcal A\).}

Indeed, for every \(S\)-scheme \(T\), there are canonical
isomorphisms
\[
  \begin{aligned}
  (\operatorname{Spec}^{*}\mathcal A^{*})(T)
  &=
  \operatorname{Hom}_{\mathcal O_T\text{-coalg.}}
    (\mathcal O_T,\mathcal A_T^{*})                                      \\
  &\simeq
  \operatorname{Hom}_{\mathcal O_T\text{-alg.}}
    (\mathcal A_T,\mathcal O_T)
  \simeq
  (\operatorname{Spec}\mathcal A)(T).
  \end{aligned}
\]

\medskip
\noindent\textbf{3.2.}
\phantomsection
\label{sga3:VIIA:3.2}
An \(\mathcal O_S\)-\emph{coalgebra in groups}, that is, a group in
the category of \(\mathcal O_S\)-coalgebras, consists of an
\(\mathcal O_S\)-coalgebra
\((\mathcal U,\Delta_{\mathcal U})\) together with three morphisms
of \(\mathcal O_S\)-coalgebras
\[
  m_{\mathcal U}:\mathcal U\otimes\mathcal U\longrightarrow\mathcal U,
  \qquad
  \eta_{\mathcal U}:\mathcal O_S\longrightarrow\mathcal U,
  \qquad
  c_{\mathcal U}:\mathcal U\longrightarrow\mathcal U,
\]
satisfying conditions \((\mathrm{ii})^{*}\),
\((\mathrm{iii})^{*}\), and~(vi) below. The requirement that
\(m_{\mathcal U}\) be a morphism of coalgebras is expressed by the
commutativity of diagrams~(iv) and~(v):
\begin{center}
  \phantomsection
  \label{sga3:VIIA:3.2-compatibility-iv-v}
  {\scriptsize
\begin{tikzcd}[column sep=small,row sep=normal]
\text{\rm(iv)}
&
\mathcal U\otimes\mathcal U
  \arrow[rr,"m_{\mathcal U}"]
  \arrow[d,"\Delta_{\mathcal U}\otimes\Delta_{\mathcal U}"']
&&
\mathcal U
  \arrow[dd,"\Delta_{\mathcal U}"]
\\
&
\mathcal U\otimes\mathcal U\otimes\mathcal U\otimes\mathcal U
  \arrow[d,"\operatorname{id}_{\mathcal U}\otimes\sigma\otimes\operatorname{id}_{\mathcal U}"']
&&
\\
&
\mathcal U\otimes\mathcal U\otimes\mathcal U\otimes\mathcal U
  \arrow[rr,"m_{\mathcal U}\otimes m_{\mathcal U}"]
&&
\mathcal U\otimes\mathcal U
\\[0.5em]
\text{\rm(v)}
&
\mathcal U\otimes\mathcal U
  \arrow[rr,"m_{\mathcal U}"]
  \arrow[dr,"\varepsilon_{\mathcal U}\otimes\varepsilon_{\mathcal U}"']
&&
\mathcal U
  \arrow[dl,"\varepsilon_{\mathcal U}"]
\\
&&
\mathcal O_S
&
\end{tikzcd}}

\end{center}

\noindent
\((\mathrm{ii})^{*}\) The square
\begin{center}
  \phantomsection
  \label{sga3:VIIA:3.2-associativity-ii-star}
  \begin{tikzcd}[column sep=huge,row sep=large]
\mathcal U\otimes\mathcal U\otimes\mathcal U
  \arrow[r,"\operatorname{id}_{\mathcal U}\otimes m_{\mathcal U}"]
  \arrow[d,"m_{\mathcal U}\otimes\operatorname{id}_{\mathcal U}"']
&
\mathcal U\otimes\mathcal U
  \arrow[d,"m_{\mathcal U}"]
\\
\mathcal U\otimes\mathcal U
  \arrow[r,"m_{\mathcal U}"]
&
\mathcal U
\end{tikzcd}

\end{center}
is commutative.

\noindent
\((\mathrm{iii})^{*}\) The two composites
\begin{center}
  \phantomsection
  \label{sga3:VIIA:3.2-unit-sequences-iii-star}
  \(
\begin{gathered}
\mathcal U\simeq\mathcal U\otimes\mathcal O_S
\xrightarrow{\operatorname{id}_{\mathcal U}\otimes\eta_{\mathcal U}}
\mathcal U\otimes\mathcal U
\xrightarrow{m_{\mathcal U}}
\mathcal U
\\[0.7em]
\mathcal U\simeq\mathcal O_S\otimes\mathcal U
\xrightarrow{\eta_{\mathcal U}\otimes\operatorname{id}_{\mathcal U}}
\mathcal U\otimes\mathcal U
\xrightarrow{m_{\mathcal U}}
\mathcal U.
\end{gathered}
\)

\end{center}
are the identity morphism of \(\mathcal U\).

\smallskip
\noindent
(vi) The composite
\begin{center}
  \phantomsection
  \label{sga3:VIIA:3.2-antipode-sequence-vi}
  \(
\mathcal U
\xrightarrow{\Delta_{\mathcal U}}
\mathcal U\otimes\mathcal U
\xrightarrow{c_{\mathcal U}\otimes\operatorname{id}_{\mathcal U}}
\mathcal U\otimes\mathcal U
\xrightarrow{m_{\mathcal U}}
\mathcal U
\,.
\)

\end{center}
is equal to \(\eta_{\mathcal U}\circ\varepsilon_{\mathcal U}\).

\medskip
\noindent\textbf{3.2.1.}
\phantomsection
\label{sga3:VIIA:3.2.1}
The morphisms \(\eta_{\mathcal U}\) and \(c_{\mathcal U}\) are
uniquely determined by \(m_{\mathcal U}\). Conditions
\((\mathrm{ii})^{*}\) and \((\mathrm{iii})^{*}\) simply say that
\(m_{\mathcal U}\) makes \(\mathcal U\) into an
\(\mathcal O_S\)-algebra whose unit section is the image under
\(\eta_{\mathcal U}\) of the unit section of \(\mathcal O_S\).
Condition~(iv) also says that the diagonal morphism
\(\Delta_{\mathcal U}\) is compatible with multiplication. Indeed,
\[
  \Delta_{\mathcal U}:\mathcal U\longrightarrow
  \mathcal U\otimes\mathcal U
\]
must be a homomorphism of coalgebras in groups, which also implies
the commutativity of the triangle
\begin{center}
  \phantomsection
  \label{sga3:VIIA:3.2.1-unit-triangle-v-star}
  \(\textup{(v)}^{*}\qquad\)
\begin{tikzcd}[column sep=huge,row sep=large,baseline=(current bounding box.center)]
&
\mathcal O_S
  \arrow[dl,"\eta_{\mathcal U}"']
  \arrow[dr,"\eta_{\mathcal U}\otimes\eta_{\mathcal U}"]
&
\\
\mathcal U
  \arrow[rr,"\Delta_{\mathcal U}"']
&& 
\mathcal U\otimes\mathcal U\,.
\end{tikzcd}

\end{center}

On the other hand, as in every category, the antipode
\(c_{\mathcal U}\) is an isomorphism from \(\mathcal U\) onto the
``opposite'' group object;\footnote{That is, the object endowed with
the multiplication
\(m'_{\mathcal U}=m_{\mathcal U}\circ\sigma\).}
in particular, \(c_{\mathcal U}\) induces an algebra isomorphism
from \(\mathcal U\) onto the opposite algebra
\(\mathcal U^{\circ}\).

\medskip
\noindent\textbf{3.2.2.}
\phantomsection
\label{sga3:VIIA:3.2.2}
Since the functor
\(\mathcal U\mapsto\operatorname{Spec}^{*}\mathcal U\) commutes
with finite products, it carries a coalgebra in groups to a
group-valued \(S\)-functor. Indeed, for every \(S\)-scheme \(T\),
the elements
\[
  x\in\Gamma(T,\mathcal U_T)
\]
belonging to
\((\operatorname{Spec}^{*}\mathcal U)(T)\) form a group under the
multiplication of the algebra \(\Gamma(T,\mathcal U_T)\), and the
inverse of \(x\) is \(c_{\mathcal U}(x)\). By \SGARefLink{sga3:VIIA:3.1.2.1}{3.1.2.1}:

\medskip
\noindent\textbf{Scholium 3.2.2.1.}
\phantomsection
\label{sga3:VIIA:3.2.2.1}
\footnote{The editors added this scholium, which is implicit in the
original.}
\textit{Let \(\mathcal U\) be an \(\mathcal O_S\)-coalgebra in
groups, finite and locally free as an \(\mathcal O_S\)-module. Then
the group-valued \(S\)-functor
\(\operatorname{Spec}^{*}\mathcal U\) is represented by the finite
locally free \(S\)-group \(\operatorname{Spec}\mathcal U^{*}\).}

\medskip
\noindent\textbf{Remark 3.2.2.2.}
\phantomsection
\label{sga3:VIIA:3.2.2.2}
Let \(\mathcal L\) be an \(\mathcal O_S\)-Lie algebra and
\(\mathcal U(\mathcal L)\) its enveloping algebra, that is, the
sheaf on \(S\) associated with the presheaf assigning to every open
subset \(V\) the enveloping algebra
\[
  U\bigl(\Gamma(V,\mathcal L)\bigr)
\]
of the Lie algebra \(\Gamma(V,\mathcal L)\).

Every homomorphism from \(\mathcal L\) to the Lie algebra underlying
an associative \(\mathcal O_S\)-algebra factors uniquely through
the canonical morphism
\(\mathcal L\to\mathcal U(\mathcal L)\). Besides the functoriality
of \(\mathcal U(\mathcal L)\) in \(\mathcal L\), this universal
property also shows that the enveloping algebra of a product of Lie
algebras is identified with the tensor product of their enveloping
algebras.

In particular, the diagonal morphism
\(\delta:\mathcal L\to\mathcal L\times\mathcal L\) induces an
algebra homomorphism
\[
  \Delta:\mathcal U(\mathcal L)\longrightarrow
  \mathcal U(\mathcal L\times\mathcal L)
  \simeq
  \mathcal U(\mathcal L)\otimes\mathcal U(\mathcal L).
\]
The zero morphism \(\mathcal L\to0\) induces a homomorphism
\[
  \varepsilon:\mathcal U(\mathcal L)\longrightarrow
  \mathcal U(0)\simeq\mathcal O_S.
\]
The isomorphism \(x\mapsto-x\) from \(\mathcal L\) onto the opposite
Lie algebra \(\mathcal L^{\circ}\) induces an anti-isomorphism \(c\)
of the algebra \(\mathcal U(\mathcal L)\). One then checks readily
that the multiplication \(m\) of \(\mathcal U(\mathcal L)\) makes
\((\mathcal U(\mathcal L),\Delta)\) into an
\(\mathcal O_S\)-coalgebra in groups with augmentation
\(\varepsilon\) and antipode \(c\).\footnote{The group-valued
\(S\)-functor \(\operatorname{Spec}^{*}\mathcal U(\mathcal L)\) is
not representable in general. Later, in \SGARefLink{sga3:VIIA:5.5}{5.5}, it will be shown that
if \(S\) is a scheme of characteristic \(p\), if \(\mathcal L\) is
finite locally free over \(\mathcal O_S\), and if
\(\mathcal U_p(\mathcal L)\) is its restricted enveloping algebra
(see \SGARefLink{sga3:VIIA:5.3}{5.3}), then
\(\operatorname{Spec}^{*}\mathcal U_p(\mathcal L)\) is represented
by a finite locally free \(S\)-group.}

\medskip
\noindent\textbf{3.2.3.}
\phantomsection
\label{sga3:VIIA:3.2.3}
\footnote{The editors expanded this paragraph.}
Let \(\mathcal U\) be an \(\mathcal O_S\)-coalgebra in groups. We
shall see that the group-valued \(S\)-functor
\[
  G=\operatorname{Spec}^{*}\mathcal U
\]
is \emph{very good}, in the sense of \SGARefLink{sga3:ii:definition:4-6}{Expos\'e II, 4.6} and \SGARefLink{sga3:ii:definition:4-10}{4.10}.

Let \(\mathcal M\) be a free \(\mathcal O_S\)-module of rank \(r\),
and let \(T\to S\) be an \(S\)-scheme. Since
\[
  I_T(\mathcal M)=
  \operatorname{Spec}(\mathcal O_T\oplus\mathcal M_T),
\]
the morphism \(\pi:I_T(\mathcal M)\to T\) is affine, and
\[
  \pi_{*}\bigl(\mathcal U_{I_T(\mathcal M)}\bigr)
  =
  \mathcal U_T\otimes_{\mathcal O_T}
  \pi_{*}\bigl(\mathcal O_{I_T(\mathcal M)}\bigr)
  =
  \mathcal U_T\otimes_{\mathcal O_T}
  (\mathcal O_T\oplus\mathcal M_T).
\]
Consequently,
\begin{equation}
  \Gamma\bigl(I_T(\mathcal M),\mathcal U_{I_T(\mathcal M)}\bigr)
  \simeq
  \Gamma(T,\mathcal U_T)\otimes_{\mathcal O(T)}
  \bigl(\mathcal O(T)\oplus\Gamma(T,\mathcal M_T)\bigr).
  \tag{1}
  \label{sga3:VIIA:3.2.3-eq1}
\end{equation}

Let \((d_1,\ldots,d_r)\) be a basis of \(\mathcal M\). An element
\[
  u_0+\sum_i u_i d_i
  \in
  \Gamma\bigl(I_T(\mathcal M),\mathcal U_{I_T(\mathcal M)}\bigr)
\]
belongs to \(G(I_T(\mathcal M))\) if and only if
\[
  1
  =\varepsilon\left(u_0+\sum_i u_i d_i\right)
  =\varepsilon(u_0)+\sum_i\varepsilon(u_i)d_i
\]
and
\[
  \left(u_0+\sum_i u_i d_i\right)
  \otimes
  \left(u_0+\sum_i u_i d_i\right)
  =
  \Delta\left(u_0+\sum_i u_i d_i\right)
  =
  \Delta(u_0)+\sum_i\Delta(u_i)d_i.
\]
Equivalently,
\begin{equation}
  \left\{
  \begin{aligned}
    \varepsilon(u_0)&=1,&
    \Delta u_0&=u_0\otimes u_0
      &&\bigl(\text{that is, }u_0\in G(T)\bigr),\\
    \varepsilon(u_i)&=0,&
    \Delta(u_i)&=u_i\otimes u_0+u_0\otimes u_i
      &&(i=1,\ldots,r).
  \end{aligned}
  \right.
  \tag{2}
  \label{sga3:VIIA:3.2.3-eq2}
\end{equation}

Moreover, the morphism
\(G(I_T(\mathcal M))\to G(T)\) corresponding to the zero section of
\(I_T(\mathcal M)\to T\) is given by
\[
  u_0+\sum_i u_i d_i\longmapsto u_0.
\]
Combining this with (1) and (2), one sees that if \(\mathcal N\) is
a second free \(\mathcal O_S\)-module of finite rank, then the
diagram of sets
\begin{center}
  \phantomsection
  \label{sga3:VIIA:3.2.3-condition-E-square}
  \begin{tikzcd}[column sep=huge,row sep=large]
G(I_T(\mathcal M\oplus\mathcal N))
  \arrow[r]
  \arrow[d]
&
G(I_T(\mathcal N))
  \arrow[d]
\\
G(I_T(\mathcal M))
  \arrow[r]
&
G(T)
\end{tikzcd}

\end{center}
is cartesian; that is, \(G\) satisfies condition~(E) of \SGARefLink{sga3:ii:definition:3-5}{Expos\'e
II, 3.5}.

Write
\[
  \operatorname{Prim}\Gamma(T,\mathcal U_T)
\]
for the \(\mathcal O(T)\)-submodule of \(\Gamma(T,\mathcal U_T)\)
formed by the \emph{primitive elements}, that is, the elements \(u\)
satisfying, with the abuse of notation noted in \SGARefLink{sga3:VIIA:3.1.2}{3.1.2},
\[
  \Delta u=u\otimes1+1\otimes u,
  \qquad
  \varepsilon(u)=0.
\]
\footnote{The second condition follows from the first, since the
first gives
\(u=(\operatorname{id}\otimes\varepsilon)\Delta(u)
=u+\varepsilon(u)\), and hence \(\varepsilon(u)=0\).}

Since \((\operatorname{Lie}G)(T)\) is the set of elements
\(u_0+ud\in G(I_T)\) above the unit element \(u_0=1\) of \(G(T)\),
there is an isomorphism of \(\mathcal O(T)\)-modules, functorial in
\(T\),\footnote{The \(\mathcal O_S\)-module structure on
\(\operatorname{Lie}G\) is defined in \SGARefLink{sga3:ii:proposition:3-6}{Expos\'e II, Proposition
3.6}.}
\[
  (\operatorname{Lie}G)(T)
  \simeq
  \operatorname{Prim}\Gamma(T,\mathcal U_T).
\]
On the other hand, \SGARefLink{sga3:VIIA:3.2.3-eq1}{(1)} yields
\[
  \operatorname{Prim}
    \Gamma\bigl(I_T(\mathcal M),
      \mathcal U_{I_T(\mathcal M)}\bigr)
  \simeq
  \operatorname{Prim}\Gamma(T,\mathcal U_T)
  \otimes_{\mathcal O(T)}\mathcal O(I_T(\mathcal M)).
\]
It follows that the natural morphism of
\(\mathcal O(I_T(\mathcal M))\)-modules
\[
  (\operatorname{Lie}G)(T)
  \otimes_{\mathcal O(T)}\mathcal O(I_T(\mathcal M))
  \longrightarrow
  (\operatorname{Lie}G)(I_T(\mathcal M))
\]
is an isomorphism; thus \(\operatorname{Lie}G\) is a good
\(\mathcal O_S\)-module (see Expos\'e II, Definition 4.4).

Therefore \(G\) is a good group-valued \(S\)-functor (see \SGARefLink{sga3:ii:definition:4-6}{Expos\'e
II, Definition 4.6}), and by \SGARefLink{sga3:ii:definition:4-7-2}{Expos\'e II, 4.7.2},
\(\operatorname{Lie}G\) carries an \(\mathcal O_S\)-bilinear Lie
bracket satisfying the Jacobi identity. It remains to show that
\(G\) is very good, that is, that the bracket on
\((\operatorname{Lie}G)(T)\) satisfies
\([u,u]=0\) for every \(u\in(\operatorname{Lie}G)(T)\); see \SGARefLink{sga3:ii:definition:4-10}{Expos\'e
II, 4.10}.

Let \(u,v\) be two elements of \((\operatorname{Lie}G)(T)\), that
is, two primitive elements of \(\Gamma(T,\mathcal U_T)\). Put
\[
  I=\operatorname{Spec}\mathcal O_S[d]/(d^2),
  \qquad
  I'=\operatorname{Spec}\mathcal O_S[d']/(d'^2).
\]
Since the composition law on \(G(I\times I')\) is induced by the
multiplication in \(\mathcal U_{I\times I'}\), one has in
\(G(I\times I')\)
\[
  \begin{aligned}
  &(1+ud)(1+vd')(1+ud)^{-1}(1+vd')^{-1}\\
  &\qquad
    =(1+ud)(1+vd')(1-ud)(1-vd')\\
  &\qquad
    =1+(uv-vu)dd'.
  \end{aligned}
\]
\begin{quote}
\small
\phantomsection
\label{sga3:VIIA:3.2.3-source-note-missing-prime}
\emph{Source note.} The re-edition prints \(1-vd\), without a prime,
in the last factor of the middle expression. Since that factor is
the inverse of \(1+vd'\), and since the displayed result is
proportional to \(dd'\), the English reconstruction makes the
required reading \(1-vd'\) explicit. Reinhold's comparison Markdown
reproduces the printed missing prime.
\end{quote}
By the description of the bracket \([u,v]\) given before
\SGARefLink{sga3:ii:proposition:4-8}{Proposition 4.8} of Expos\'e II, it follows that
\[
  [u,v]=uv-vu,
\]
where the right-hand side is the commutator of \(u\) and \(v\) in
the algebra \(\Gamma(T,\mathcal U_T)\). Hence \([u,u]=0\), and we
have proved:

\medskip
\noindent\textbf{Proposition.}
\phantomsection
\label{sga3:VIIA:3.2.3-primitive-elements-proposition}
\textit{Let \(\mathcal U\) be an \(\mathcal O_S\)-coalgebra in
groups. The group-valued \(S\)-functor
\(G=\operatorname{Spec}^{*}\mathcal U\) is very good, and there is
an isomorphism
\[
  \operatorname{Lie}G
  \simeq
  \operatorname{Prim}\mathcal W(\mathcal U)
\]
of \(\mathcal O_S\)-Lie algebras, where
\(\operatorname{Prim}\mathcal W(\mathcal U)\) denotes the functor
which assigns to every \(T\to S\) the \(\mathcal O(T)\)-Lie algebra
of primitive elements of
\(\mathcal W(\mathcal U)(T)=\Gamma(T,\mathcal U_T)\).}
\footnote{The editors added this proposition to summarize the
preceding discussion.}

\medskip
\noindent\textbf{3.3.}
\phantomsection
\label{sga3:VIIA:3.3}
Suppose finally that \(\mathcal U\) is a
\emph{coalgebra in commutative groups}, that is, that the triangle
\begin{center}
  \phantomsection
  \label{sga3:VIIA:3.3-commutative-triangle-i-star}
  \(\textup{(i)}^{*}\qquad\)
\begin{tikzcd}[column sep=huge,row sep=large,baseline=(current bounding box.center)]
\mathcal U\otimes\mathcal U
  \arrow[rr,"\sigma"]
  \arrow[dr,"m_{\mathcal U}"']
&&
\mathcal U\otimes\mathcal U
  \arrow[dl,"m_{\mathcal U}"]
\\
&
\mathcal U
&
\end{tikzcd}

\end{center}
is commutative, or equivalently that \(m_{\mathcal U}\) makes
\(\mathcal U\) into a commutative \(\mathcal O_S\)-algebra.
Conditions \SGARefLink{sga3:vii:paragraph:5-1:equation:i}{(i)}, \SGARefLink{sga3:vii:paragraph:5-1:equation:ii}{(ii)}, (iii), (iv), (v), (vi),
\((\mathrm{i})^{*}\), \((\mathrm{ii})^{*}\),
\((\mathrm{iii})^{*}\), and \((\mathrm{v})^{*}\) then also say that
\(\mathcal U\) is a cogroup in the category of commutative
\(\mathcal O_S\)-algebras. Thus, if in addition \(\mathcal U\) is a
quasi-coherent \(\mathcal O_S\)-module, the affine \(S\)-scheme
\(\operatorname{Spec}\mathcal U\) is a commutative \(S\)-group
scheme.

In this case the diagonal morphism \(\Delta'\) of
\(\mathcal O_S[T,T^{-1}]\) sends \(T\) to \(T\otimes T\).
Consequently, morphisms of \(S\)-groups from
\(\operatorname{Spec}\mathcal U\) to
\(\mathbb G_{m,S}\) (Expos\'e I, 4.3.2) correspond bijectively to
morphisms of unital \(\mathcal O_S\)-algebras
\[
  \varphi:\mathcal O_S[T,T^{-1}]\longrightarrow\mathcal U
\]
such that
\[
  (\varphi\otimes\varphi)\circ\Delta'
  =\Delta_{\mathcal U}\circ\varphi.
\]
In this situation,
\(\varepsilon_{\mathcal U}\circ\varphi\) is the neutral element of
\(\mathbb G_{m,S}(S)\), that is, the augmentation
\(\varepsilon'\). Such a morphism \(\varphi\) is determined by the
image \(\varphi(T)\), which must be an invertible element
\(x\in\mathcal U\) satisfying
\[
  \Delta_{\mathcal U}x=x\otimes x,
  \qquad
  \varepsilon_{\mathcal U}(x)=\varepsilon'(T)=1.
\]
Therefore
\[
  \operatorname{Hom}_{S\text{-gr.}}
  (\operatorname{Spec}\mathcal U,\mathbb G_{m,S})
  \simeq
  (\operatorname{Spec}^{*}\mathcal U)(S).
\]
Since this formula remains valid after arbitrary base change, it
gives
\[
  \operatorname{Spec}^{*}\mathcal U
  \simeq
  \operatorname{Hom}_{S\text{-gr.}}
  (\operatorname{Spec}\mathcal U,\mathbb G_{m,S}).
\]
We have therefore proved:

\medskip
\noindent\textbf{Proposition 3.3.0.}
\phantomsection
\label{sga3:VIIA:3.3.0}
\textit{If \(\mathcal U\) is an \(\mathcal O_S\)-coalgebra in
commutative groups, quasi-coherent as an \(\mathcal O_S\)-module,
then the affine \(S\)-scheme
\(G=\operatorname{Spec}\mathcal U\) is a commutative \(S\)-group
scheme, and there is an isomorphism of group-valued \(S\)-functors
\[
  \operatorname{Spec}^{*}\mathcal U
  \simeq
  \operatorname{Hom}_{S\text{-gr.}}(G,\mathbb G_{m,S}).
\]}

If, moreover, \(\mathcal U\) is locally free of finite type as an
\(\mathcal O_S\)-module, then by \SGARefLink{sga3:VIIA:3.1.2.1}{3.1.2.1} the group-valued
\(S\)-functor \(\operatorname{Spec}^{*}\mathcal U\) is represented
by \(\operatorname{Spec}\mathcal U^{*}\). Hence:

\medskip
\noindent\textbf{Proposition 3.3.1 (Cartier duality).}
\phantomsection
\label{sga3:VIIA:3.3.1}
\textit{The functor
\[
  \mathscr A(G)\longmapsto
  \mathscr A(G)^{*}
  =
  \mathcal Hom_{\mathcal O_S\text{-Mod.}}
  (\mathscr A(G),\mathcal O_S)
\]
induces a duality\textsuperscript{\(\ast\)} on the category of
commutative finite locally free \(S\)-group schemes; it associates
to \(G\) the \(S\)-group
\(\operatorname{Hom}_{S\text{-gr.}}(G,\mathbb G_{m,S})\).}

\begingroup
\renewcommand{\thefootnote}{\fnsymbol{footnote}}
\footnotetext[1]{A duality of a category \(\mathscr C\) is a pair
\((D,\varphi)\) consisting of a contravariant functor
\(D:\mathscr C\to\mathscr C\) and a natural isomorphism
\(\varphi:\operatorname{Id}_{\mathscr C}\to DD\), such that the
isomorphisms
\(\varphi D:D\to DDD\) and
\(D\varphi^{-1}:DDD\to D\) are inverse to one another. The editors
corrected \(D\varphi\) to \(D\varphi^{-1}\).}
\endgroup

